\section{第一届}

\subsection*{一、填空题(本题 20分,  每小题 5分,  共 4 小题)}

\begin{enumerate}
    \item 计算 $\displaystyle\iint\limits_D \frac{(x + y)\ln\left(1 + \frac{y}{x}\right)}{\sqrt{1 - x - y}}  \mathrm{d}x\mathrm{d}y = \underline{\quad\quad}$,  其中区域 $D$ 由直线 $x + y = 1$ 与两坐标轴所围三角形区域.
    \item 设 $f(x)$ 是连续函数,  满足
          $f(x) = 3x^2 - \displaystyle\int_0^2 f(x)\mathrm{d}x - 2, $
          则 $f(x) = \underline{\quad\quad}$.
    \item 曲面 $z = \dfrac{x^2}{2} + y^2 - 2$ 平行平面 $2x + 2y - z = 0$ 的切平面方程是 $\underline{\quad\quad}$.
    \item 设 $y = y(x)$ 由方程 $xe^{f(y)} = e^y \ln 29$ 确定,  其中 $f$ 具有二阶导数,  且 $f' \neq 1$,  则 $\dfrac{\mathrm{d}^2 y}{\mathrm{d}x^2} = \underline{\quad\quad}$.
\end{enumerate}

\subsection*{二、(本题5分)}

\begin{enumerate}[start=5]
    \item 求极限 $\displaystyle\lim_{x \to 0} \left( \dfrac{e^x + e^{2x} + \cdots + e^{nx}}{n} \right)^{\frac{e}{x}}$,  其中 $n$ 是给定的正整数.
\end{enumerate}

\subsection*{三、(本题15分)}

\begin{enumerate}[start=6]
    \item 设函数 $f(x)$ 连续,  $g(x) = \displaystyle\int_0^1 f(xt)  \mathrm{d}t$,  且 $\displaystyle\lim_{x \to 0} \dfrac{f(x)}{x} = A$,  $A$ 为常数,  求 $g'(x)$ 并讨论 $g'(x)$ 在 $x = 0$ 处的连续性.
\end{enumerate}

\subsection*{四、(本题15分)}

\begin{enumerate}[start=7]
    \item 已知平面区域
          $D = \{(x,  y) \mid 0 \leq x \leq \pi,  0 \leq y \leq \pi\},  $
          $L$ 为 $D$ 的正向边界,  试证：
          \begin{enumerate}
              \item [(1)] $\displaystyle\oint_L x e^{\sin y}  \mathrm{d}y - y e^{-\sin x}  \mathrm{d}x = \oint_L x e^{-\sin y}  \mathrm{d}y - y e^{\sin x}  \mathrm{d}x$.
              \item [(2)] $\displaystyle\oint_L x e^{\sin y}  \mathrm{d}y - y e^{-\sin x}  \mathrm{d}x \geq \frac{5}{2} \pi^2$.
          \end{enumerate}
\end{enumerate}

\subsection*{五、(本题10分)}

\begin{enumerate}[start=8]
    \item 已知 $y_1 = x e^x + e^{2x}$,  $y_2 = x e^x + e^{-x}$,  $y_3 = x e^x + e^{2x} - e^{-x}$ 是某二阶常系数线性非齐次微分方程的三个解,  试求此微分方程.
\end{enumerate}

\subsection*{六、(本题10分)}

\begin{enumerate}[start=9]
    \item 设抛物线 $y = ax^2 + bx + 2 \ln c$ 过原点,  当 $0 \leq x \leq 1$ 时,  $y \geq 0$,  又已知该抛物线与 $x$ 轴及直线 $x = 1$ 所围图形的面积为 $\dfrac{1}{3}$. 试确定 $a,  b,  c$ 使此图形绕 $x$ 轴旋转一周而成的旋转体的体积 $V$ 最小.
\end{enumerate}

\subsection*{七、(本题15分)}

\begin{enumerate}[start=10]
    \item 已知 $u_n(x)$ 满足 $u_n'(x) = u_n(x) + x^{n - 1} e^x$($n$ 为正整数),  且 $u_n(1) = \dfrac{e}{n}$,  求函数项级数 $\displaystyle\sum_{n = 1}^{\infty} u_n(x)$ 之和.
\end{enumerate}

\subsection*{八、(本题10分)}

\begin{enumerate}[start=11]
    \item 求 $x \to 1^-$ 时,  与 $\displaystyle\sum_{n = 0}^{\infty} x^{n^2}$ 等价的无穷大量.
\end{enumerate}

\newpage
\section{第二届}

\subsection*{一、计算题(本题 25 分,  每小题5分,  共5小题)}

\begin{enumerate}
    \item 设 $x_n = (1 + a) \cdot (1 + a^2) \cdots (1 + a^{2^n})$,  其中 $|a| < 1$,  求 $\displaystyle\lim_{n \to \infty} x_n$.
    \item 求 $\displaystyle\lim_{x \to \infty} e^{-x} \left( 1 + \dfrac{1}{x} \right)^{x^2}$.
    \item 设 $s > 0$,  求 $I_n = \displaystyle\int_0^{+\infty} e^{-sx} x^n  \mathrm{d}x \ (n = 1,  2,  \cdots)$.
    \item 设 $f(t)$ 有二阶连续导数,  $r = \sqrt{x^2 + y^2}$,  $g(x,  y) = f\left( \dfrac{1}{r} \right)$,  求 $\dfrac{\partial^2 g}{\partial x^2} + \dfrac{\partial^2 g}{\partial y^2}$.
    \item 求直线 $l_1: \begin{cases} x - y = 0 \\ z = 0 \end{cases}$ 与直线 $l_2: \dfrac{x - 2}{4} = \dfrac{y - 1}{-2} = \dfrac{z - 3}{-1}$ 的距离.
\end{enumerate}

\subsection*{二、(本题15分)}

\begin{enumerate}[start=6]
    \item 设函数 $f(x)$ 在 $(-\infty,  +\infty)$ 上具有二阶导数,  并且
          \[
              f''(x) > 0, \ \displaystyle\lim_{x \to +\infty} f'(x) = \alpha > 0, \ \displaystyle\lim_{x \to -\infty} f'(x) = \beta < 0,
          \]
          且存在一点 $x_0$,  使得 $f(x_0) < 0$. 证明：方程 $f(x) = 0$ 在 $(-\infty,  +\infty)$ 恰有两个实根.
\end{enumerate}

\subsection*{三、(本题15分)}

\begin{enumerate}[start=7]
    \item 设 $y = f(x)$ 由参数方程
          $\begin{cases}
                  x = 2t + t^2 , \\
                  y = \psi(t)
              \end{cases}  (t > -1)$
          所确定. 且 $\dfrac{\mathrm{d}^2 y}{\mathrm{d}x^2} = \dfrac{3}{4(1 + t)}$,  其中 $\psi(t)$ 具有二阶导数,  曲线 $y = \psi(t)$ 与 $y = \displaystyle\int_1^{t^2} e^{-u^2}  \mathrm{d}u + \frac{3}{2e}$ 在 $t = 1$ 处相切. 求函数 $\psi(t)$.
\end{enumerate}

\subsection*{四、(本题15分)}

\begin{enumerate}[start=8]
    \item 设 $a_n > 0$,  $S_n = \displaystyle\sum_{k = 1}^n a_k$,  证明：
          \begin{enumerate}
              \item [(1)] 当 $\alpha > 1$ 时,  级数 $\displaystyle\sum_{n = 1}^{+\infty} \dfrac{a_n}{S_n^\alpha}$ 收敛.
              \item [(2)] 当 $\alpha \leq 1$,  且 $S_n \to \infty \ (n \to \infty)$ 时,  $\displaystyle\sum_{n = 1}^{+\infty} \dfrac{a_n}{S_n^\alpha}$ 发散.
          \end{enumerate}
\end{enumerate}

\subsection*{五、(本题15分)}

\begin{enumerate}[start=9]
    \item 设 $l$ 是过原点、方向为 $(\alpha,  \beta,  \gamma)$(其中 $\alpha^2 + \beta^2 + \gamma^2 = 1$)的直线,  均匀椭球 $\dfrac{x^2}{a^2} + \dfrac{y^2}{b^2} + \dfrac{z^2}{c^2} \leq 1$(其中 $0 < c < b < a$,  密度为 1)绕 $l$ 旋转.
          \begin{enumerate}
              \item [(1)] 求其转动惯量.
              \item [(2)] 求其转动惯量关于方向 $(\alpha,  \beta,  \gamma)$ 的最大值和最小值.
          \end{enumerate}
\end{enumerate}

\subsection*{六、(本题15分)}

\begin{enumerate}[start=10]
    \item 设函数 $\varphi(x)$ 具有连续的导数,  在围绕原点的任意光滑的简单闭曲线 $C$ 上,  积分 $\displaystyle\oint_C \frac{2xy  \mathrm{d}x + \varphi(x)  \mathrm{d}y}{x^4 + y^2}$ 的值为常数.
          \begin{enumerate}
              \item [(1)]设 $L$ 为正向闭曲线 $(x - 2)^2 + y^2 = 1$. 证明：
                    $\displaystyle\oint_L \frac{2xy  \mathrm{d}x + \varphi(x)  \mathrm{d}y}{x^4 + y^2} = 0.$
              \item [(2)] 求函数 $\varphi(x)$.
              \item [(3)] 设 $C$ 是围绕原点的光滑简单正向闭曲线,  求 $\displaystyle\oint_C \frac{2xy  \mathrm{d}x + \varphi(x)  \mathrm{d}y}{x^4 + y^2}$.
          \end{enumerate}
\end{enumerate}


\newpage
\section{第三届}

\subsection*{一、计算题(本题 24 分,  每小题 6分,  共 4小题)}
\begin{enumerate}
    \item $\displaystyle\lim_{x \to 0} \frac{(1+x)^{\frac{2}{x}} - e^2(1 - \ln(1+x))}{x}.$
    \item 设 $a_n = \cos \dfrac{\theta}{2} \cdot \cos \dfrac{\theta}{2^2} \cdots \cos \dfrac{\theta}{2^n}$,  求 $\displaystyle\lim_{n \to \infty} a_n$.
    \item 求
          $\displaystyle\iint\limits_D \text{sgn}(xy-1) \mathrm{d}x  \mathrm{d}y, $
          其中 $D = \{(x,  y) \mid 0 \leq x \leq 2,  0 \leq y \leq 2\}$.
    \item 求幂级数
          $\displaystyle\sum_{n=1}^{\infty} \dfrac{2n-1}{2^n} x^{2n-2}$ 的和函数,  并求级数
          $\displaystyle\sum_{n=1}^{\infty} \dfrac{2n-1}{2^{2n-1}}$ 的和.
\end{enumerate}

\subsection*{二、(本题 16 分,  每小题8 分)}
\begin{enumerate}[start=5]
    \item 设 $\{a_n\}_{n=0}^{\infty}$ 为数列,  $a,  \lambda$ 为有限数,  求证：
          \begin{enumerate}
              \item [(1)] 如果 $\displaystyle\lim_{n \to \infty} a_n = a$,  则 $\displaystyle\lim_{n \to \infty} \frac{a_1 + a_2 + \cdots + a_n}{n} = a$.
              \item [(2)] 如果存在正整数 $p$,  使得 $\displaystyle\lim_{n \to \infty} (a_{n+p} - a_n) = \lambda$,  则 $\displaystyle\lim_{n \to \infty} \dfrac{a_n}{n} = \dfrac{\lambda}{p}$.
          \end{enumerate}
\end{enumerate}

\subsection*{三、(本题15分)}
\begin{enumerate}[start=6]
    \item 设函数 $f(x)$ 在闭区间 $[-1, 1]$ 上具有连续的三阶导数,  且 $f(-1) = 0,  f(1) = 1,  f'(0) = 0$.求证：在开区间 $(-1, 1)$ 内至少存在一点 $x_0$,  使得 $f'''(x_0) = 3$.
\end{enumerate}

\subsection*{四、(本题15分)}
\begin{enumerate}[start=7]
    \item 在 $xOy$ 平面上,  有一条从点 $(a, 0)$ 向右的射线,  线密度为 $\rho$.在点 $(0, h)$ 处(其中 $h > 0$)有一质量为 $m$ 的质点.求射线对该质点的引力.
\end{enumerate}

\subsection*{五、(本题15分)}
\begin{enumerate}[start=8]
    \item 设 $z = z(x,  y)$ 是由方程 $F\left(z + \dfrac{1}{x},  z - \dfrac{1}{y}\right) = 0$ 确定的隐函数,  且具有连续的二阶偏导数.求证：$x^2 \dfrac{\partial z}{\partial x} - y^2 \dfrac{\partial z}{\partial y} = 1$ 和 $x^3 \dfrac{\partial^2 z}{\partial x^2} + xy(x-y) \dfrac{\partial^2 z}{\partial x \partial y} - y^3 \dfrac{\partial^2 z}{\partial y^2} = -2$.
\end{enumerate}

\subsection*{六、(本题15分)}
\begin{enumerate}[start=9]
    \item 设函数 $f(x)$ 连续,  $a,  b,  c$ 为常数,  $\Sigma$ 是单位球面 $x^2 + y^2 + z^2 = 1$.求证：
          \[
              \iint\limits_{\Sigma} f(ax + by + cz)  \mathrm{d}S = 2\pi \int_{-1}^{1} f\left(\sqrt{a^2 + b^2 + c^2} u\right)  \mathrm{d}u.
          \]
\end{enumerate}


\newpage
\section{第四届}

\subsection*{一、解答题(本题30分,  每小题6分,  共5小题)}
\begin{enumerate}
    \item 求极限 $\displaystyle\lim_{n \to \infty} (n!)^{\frac{1}{n^2}}$.
    \item 求通过直线 $L:
              \begin{cases}
                  2x + y - 3z + 2 = 0, \\
                  5x + 5y - 4z + 3 = 0
              \end{cases}$ 的两个相互垂直的平面 $\pi_1$ 和 $\pi_2$,  使其中一个平面过点 $(4,  -3,  1)$.
    \item 已知函数 $z = u(x,  y)e^{ax + by}$,  且 $\dfrac{\partial^2 u}{\partial x \partial y} = 0$,  确定常数 $a$ 和 $b$,  使函数 $z = z(x,  y)$ 满足方程
          $\dfrac{\partial^2 z}{\partial x \partial y} - \dfrac{\partial z}{\partial x} - \dfrac{\partial z}{\partial y} + z = 0.$
    \item 设函数 $u = u(x)$ 连续可微,  $u(2) = 1$,  且 $\displaystyle\int_L (x + 2y) u  \mathrm{d}x + (x + u^3) u  \mathrm{d}y$ 在右半平面上与路径无关,  求 $u(x)$.
    \item 求极限 $\displaystyle\lim_{x \to +\infty} \sqrt[3]{x} \int_x^{x+1} \frac{\sin t}{\sqrt{t + \cos t}}  \mathrm{d}t$.
\end{enumerate}

\subsection*{二、(本题10分)}
\begin{enumerate}[start=6]
    \item 计算 $\displaystyle\int_0^{+\infty} e^{-2x} |\sin x|  \mathrm{d}x$.
\end{enumerate}

\subsection*{三、(本题10分)}
\begin{enumerate}[start=7]
    \item 求方程 $x^2 \sin \dfrac{1}{2} = 2x - 501$ 的近似解,  精确到 0.001.
\end{enumerate}

\subsection*{四、(本题12分)}
\begin{enumerate}[start=8]
    \item 设函数 $y = f(x)$ 二阶可导,  且 $f''(x) > 0$,  又设 $f(0) = 0,  f'(0) = 0$.求 $\displaystyle\lim_{x \to 0} \frac{x^3 f(u)}{f(x) \sin^3 u}$,  其中 $u$ 是曲线 $y = f(x)$ 上点 $P(x,  f(x))$ 处的切线在 $x$ 轴上的截距.
\end{enumerate}

\subsection*{五、(本题12分)}
\begin{enumerate}[start=9]
    \item 求最小实数 $C$,  使得满足 $\displaystyle\int_0^1 |f(x)|  \mathrm{d}x = 1$ 的连续函数 $f(x)$ 都有
          \[
              \int_0^1 f(\sqrt{x})  \mathrm{d}x \leqslant C.
          \]
\end{enumerate}

\subsection*{六、(本题12分)}
\begin{enumerate}[start=10]
    \item 设 $f(x)$ 为连续函数,  $t > 0$.区域 $\Omega$ 由抛物面 $z = x^2 + y^2$ 和球面 $x^2 + y^2 + z^2 = t^2$ 所围成.定义三重积分
          \[
              F(t) = \iiint\limits_{\Omega} f(x^2 + y^2 + z^2)  \mathrm{d}v,
          \]
          求 $F(t)$ 的导数 $F'(t)$.
\end{enumerate}

\subsection*{七、(本题14分)}
\begin{enumerate}[start=11]
    \item 设 $\displaystyle\sum_{n=1}^{\infty} a_n$ 与 $\displaystyle\sum_{n=1}^{\infty} b_n$ 为正项级数.
          \begin{enumerate}
              \item [(1)] 若 $\displaystyle\lim_{n \to \infty} \left( \frac{a_n}{a_{n+1}b_n} - \frac{1}{b_{n+1}} \right) > 0$,  则 $\displaystyle\sum_{n=1}^{\infty} a_n$ 收敛.
              \item [(2)] 若 $\displaystyle\lim_{n \to \infty} \left( \frac{a_n}{a_{n+1}b_n} - \frac{1}{b_{n+1}} \right) < 0$,  且 $\displaystyle\sum_{n=1}^{\infty} b_n$ 发散,  则 $\displaystyle\sum_{n=1}^{\infty} a_n$ 发散.
          \end{enumerate}
\end{enumerate}


\newpage
\section{第五届}

\subsection*{一、解答题(本题24分,  每小题6分,  共4小题)}
\begin{enumerate}
    \item 求极限 $\displaystyle\lim_{n \to \infty} \left(1 + \sin(\pi \sqrt{1 + 4n^2})\right)^n$.
    \item 证明广义积分 $\displaystyle\int_0^{+\infty} \frac{\sin x}{x}  \mathrm{d}x$ 不是绝对收敛的.
    \item 设函数 $y = y(x)$ 由 $x^3 + 3x^2y - 2y^3 = 2$ 所确定,  求 $y(x)$ 的极值.
    \item 过曲线 $y = \sqrt[3]{x} (x \geq 0)$ 上的点 $A$ 作切线,  使该切线与曲线及 $x$ 轴所围成的平面图形的面积为 $\dfrac{3}{4}$,  求点 $A$ 的坐标.
\end{enumerate}

\subsection*{二、(本题 12 分)}
\begin{enumerate}[start=5]
    \item 计算定积分 $I = \displaystyle\int_{-\pi}^{\pi} \frac{x \sin x \cdot \arctan e^x}{1 + \cos^2 x}  \mathrm{d}x$.
\end{enumerate}

\subsection*{三、(本题 12 分)}
\begin{enumerate}[start=6]
    \item 设 $f(x)$ 在 $x = 0$ 处存在二阶导数,  且 $\displaystyle\lim_{x \to 0} \frac{f(x)}{x} = 0$.证明：级数 $\displaystyle\sum_{n=1}^{\infty} \left| f\left( \frac{1}{n} \right) \right|$ 收敛.
\end{enumerate}

\subsection*{四、(本题 10 分)}
\begin{enumerate}[start=7]
    \item 设 $|f(x)| \leqslant \pi,  f'(x) \geqslant m > 0  (\forall x \in [a,  b])$,  证明：
          \[
              \left| \int_a^b \sin f(x)  \mathrm{d}x \right| \leqslant \frac{2}{m}.
          \]
\end{enumerate}

\subsection*{五、(本题 14 分)}
\begin{enumerate}[start=8]
    \item 设 $\Sigma$ 是一个光滑封闭曲面,  方向朝外,  并给定第二型曲面积分
          \[
              I = \iint\limits_{\Sigma} (x^3 - x)  \mathrm{d}y  \mathrm{d}z + (2y^3 - y)  \mathrm{d}z  \mathrm{d}x + (3z^3 - z)  \mathrm{d}x  \mathrm{d}y.
          \]
          试确定曲面 $\Sigma$,  使得积分 $I$ 的值最小,  并求该最小值.
\end{enumerate}

\subsection*{六、(本题 14 分)}
\begin{enumerate}[start=9]
    \item 设 $I_a(r) = \displaystyle\int_C \frac{y  \mathrm{d}x - x  \mathrm{d}y}{(x^2 + y^2)^\alpha}$,  其中 $a$ 为常数,  曲线 $C$ 为椭圆 $x^2 + xy + y^2 = r^2$,  取正向.求极限 $\displaystyle\lim_{r \to +\infty} I_a(r)$.
\end{enumerate}

\subsection*{七、(本题 14 分)}
\begin{enumerate}[start=10]
    \item 判断级数 $\displaystyle\sum_{n=1}^{\infty} \frac{1}{(n+1)(n+2)} + \frac{1}{2} + \cdots + \frac{1}{n}$ 的敛散性,  若收敛,  则求其和.
\end{enumerate}


\newpage
\section{第六届}

\subsection*{一、填空题(本题30 分,  每小题 6分,  共 5小题)}
\begin{enumerate}
    \item 已知 $y_1 = e^x$ 和 $y_2 = xe^x$ 是齐次二阶常系数线性微分方程的解,  则该方程为
    \item 设有曲面 $S: z = x^2 + 2y^2$ 和平面 $L: 2x + 2y + z = 0$,  则与 $L$ 平行的 $S$ 的切平面是
    \item 设函数 $y = y(x)$ 由方程 $x = \displaystyle\int_1^{y-x} \sin^2 \frac{\pi t}{4}  \mathrm{d}t$ 所确定,  求 $\left. \dfrac{\mathrm{d}y}{\mathrm{d}x} \right|_{x=0} = \underline{\quad\quad}$.
    \item 设 $x_n = \displaystyle\sum_{k=1}^{n} \frac{k}{(k+1)!}$,  则 $\displaystyle\lim_{n \to \infty} x_n = \underline{\quad\quad}$.
    \item 已知 $\displaystyle\lim_{x \to 0} \left(1 + x + \frac{f(x)}{x}\right)^{\frac{1}{x}} = e^3$,  则 $\displaystyle\lim_{x \to 0} \frac{f(x)}{x^2} = \underline{\quad\quad}$.
\end{enumerate}

\subsection*{二、(本题12分)}
\begin{enumerate}[start=6]
    \item 设 $n$ 为正整数,  计算 $I = \displaystyle\int_{e^{-2n\pi}}^1 \left| \frac{\mathrm{d}}{\mathrm{d}x} \cos \left( \ln \frac{1}{x} \right) \right| \mathrm{d}x$.
\end{enumerate}

\subsection*{三、(本题14分)}
\begin{enumerate}[start=7]
    \item 设函数 $f(x)$ 在 $[0, 1]$ 上有二阶导数,  且存在正常数 $A,  B$ 使得 $|f(x)| \leqslant A,  |f''(x)| \leqslant B$.证明：对任意 $x \in [0, 1]$,  有 $|f'(x)| \leqslant 2A + \dfrac{B}{2}$.
\end{enumerate}

\subsection*{四、(本题14分)}
\begin{enumerate}[start=8]
    \item
          \begin{enumerate}
              \item [(1)] 设一球缺高为 $h$,  所在球半径为 $R$.证明该球缺的体积为 $\dfrac{\pi}{3}(3R-h)h^2$,  球冠的面积为 $2\pi Rh$.
              \item [(2)] 设球体 $(x-1)^2 + (y-1)^2 + (z-1)^2 \leqslant 12$ 被平面 $P: x + y + z = 6$ 所截的小球缺为 $\Omega$.记球上缺的球冠为 $\Sigma$,  方向指向球外,  求第二型曲面积分
                    \[
                        I = \iint\limits_{\Sigma} x  \mathrm{d}y  \mathrm{d}z + y  \mathrm{d}z  \mathrm{d}x + z  \mathrm{d}x  \mathrm{d}y.
                    \]
          \end{enumerate}
\end{enumerate}

\subsection*{五、(本题15分)}
\begin{enumerate}[start=9]
    \item 设 $f$ 在 $[a, b]$ 上非负连续,  严格单增,  且对于任意的 $n \in \mathbb{N}$,  存在 $x_n \in [a, b]$ 使得
          \[
              [f(x_n)]^n = \frac{1}{b-a} \int_a^b [f(x)]^n  \mathrm{d}x.
          \]
          求 $\displaystyle\lim_{n \to \infty} x_n$.
\end{enumerate}

\subsection*{六、(本题15分)}
\begin{enumerate}[start=10]
    \item 设 $A_n = \dfrac{n}{n^2 + 1} + \dfrac{n}{n^2 + 2^2} + \cdots + \dfrac{n}{n^2 + n^2}$,  求 $\displaystyle\lim_{n \to \infty} n \left( \frac{\pi}{4} - A_n \right)$.
\end{enumerate}


\newpage
\section{第七届}

\subsection*{一、填空题(本题 30 分,  每小题 6分,  共5 小题)}
\begin{enumerate}
    \item (1) 求极限 $\displaystyle\lim_{n \to \infty} n \left( \frac{\sin \dfrac{\pi}{n}}{n^2 + 1} + \frac{\sin \dfrac{2\pi}{n}}{n^2 + 2} + \cdots + \frac{\sin \pi}{n^2 + n} \right) = \underline{\quad\quad}. $
    \item 设函数 $z = z(x,  y)$ 由方程 $F\left(x + \dfrac{z}{y},  y + \dfrac{z}{x}\right) = 0$ 所确定,  其中 $F(u,  v)$ 具有连续偏导数且 $xF_u + yF_v \neq 0$,  则 $x \dfrac{\partial z}{\partial x} + y \dfrac{\partial z}{\partial y} = \underline{\quad\quad}$(本小题结果要求不显含 $F$ 及其偏导数).
    \item 曲面 $z = x^2 + y^2 + 1$ 在点 $M(1,  -1,  3)$ 处的切平面与曲面 $z = x^2 + y^2$ 所围区域的体积为 \underline{\quad\quad}.
    \item 函数 $f(x) =
              \begin{cases}
                  3, & x \in [-5,  0), \\
                  0, & x \in [0,  5)
              \end{cases}$ 在 $(-5,  5]$ 上的 Fourier 级数在 $x = 0$ 处收敛的值是 \underline{\quad\quad}.
    \item 设区间 $(0,  +\infty)$ 上的函数 $u(x)$ 定义为 $u(x) = \displaystyle\int_0^{+\infty} e^{-xt^2}  \mathrm{d}t$,  则 $u(x)$ 的初等表达式是 \underline{\quad\quad}.

\end{enumerate}

\subsection*{二、(本题12分)}
\begin{enumerate}[start=6]
    \item 设 $M$ 是以三个正半坐标轴为母线的半圆锥面,  求其方程.
\end{enumerate}

\subsection*{三、(本题12分)}
\begin{enumerate}[start=7]
    \item 设 $f(x)$ 在 $(a,  b)$ 内二阶可导,  且存在常数 $\alpha,  \beta$,  使得对于 $\forall x \in (a,  b)$,  有 $f'(x) = \alpha f(x) + \beta f''(x)$.证明：$f(x)$ 在 $(a,  b)$ 内无穷阶可导.
\end{enumerate}

\subsection*{四、(本题14分)}
\begin{enumerate}[start=8]
    \item 求幂级数 $\displaystyle\sum_{n=0}^{\infty} \frac{n^3 + 2}{(n + 1)!} (x - 1)^n$ 的收敛域及其和函数.
\end{enumerate}

\subsection*{五、(本题16分)}
\begin{enumerate}[start=9]
    \item 设函数 $f$ 在 $[0,  1]$ 上连续,  且 $\displaystyle\int_0^1 f(x)  \mathrm{d}x = 0,  \displaystyle\int_0^1 x f(x)  \mathrm{d}x = 1$.证明：
          \begin{enumerate}
              \item [(1)] 存在 $x_0 \in [0,  1]$,  使得 $|f(x_0)| > 4$.
              \item [(2)] 存在 $x_1 \in [0,  1]$,  使得 $|f(x_1)| = 4$.
          \end{enumerate}
\end{enumerate}

\subsection*{六、(本题16分)}
\begin{enumerate}[start=10]
    \item 设 $f(x,  y)$ 在 $x^2 + y^2 \leqslant 1$ 上有连续的二阶偏导数,  且 $f_{xx}^2 + 2f_{xy}^2 + f_{yy}^2 \leqslant M$.若 $f(0,  0) = 0,  f_x(0,  0) = 0,  f_y(0,  0) = 0$,  证明：
          \[
              \left| \iint\limits_{x^2 + y^2 \leqslant 1} f(x,  y)  \mathrm{d}x  \mathrm{d}y \right| \leqslant \frac{\pi \sqrt{M}}{4}.
          \]
\end{enumerate}


\newpage
\section{第八届}

\subsection*{一、填空题(本题 30 分,  每小题 6分,  共 5小题)}
\begin{enumerate}
    \item 若 $f(x)$ 在点 $x = a$ 处可导,  且 $f(a) \neq 0$,  则 $\displaystyle\lim_{n \to \infty} \left( \frac{f\left(a + \frac{1}{n}\right)}{f(a)} \right)^n = \underline{\quad\quad}$.
    \item 若 $f(1) = 0,  f'(1)$ 存在,  则极限 $\displaystyle\lim_{x \to 0} \frac{f\left(\sin^2 x + \cos x\right) \tan 3x}{(e^{x^2} - 1) \sin x} = \underline{\quad\quad}$.
    \item 设 $f(x)$ 有连续导数,  且 $f(1) = 2$.记 $z = f\left(e^x y^2\right)$.若 $\dfrac{\partial z}{\partial x} = z$,  则当 $x > 0$ 时,  $f(x) = \underline{\quad\quad}$.
    \item 设 $f(x) = e^x \sin 2x$,  则 $f^{(4)}(0) = \underline{\quad\quad}$.
    \item 曲面 $z = \dfrac{x^2}{2} + y^2$ 平行于平面 $2x + 2y - z = 0$ 的切平面方程为 \underline{\quad\quad}.
\end{enumerate}

\subsection*{二、(本题14分)}
\begin{enumerate}[start=6]
    \item 设 $f(x)$ 在 $[0, 1]$ 上可导,  $f(0) = 0$,  且当 $x \in (0, 1),  0 < f'(x) < 1$.试证：当 $a \in (0, 1)$,  $\left( \displaystyle\int_0^a f(x)  \mathrm{d}x \right)^2 > \displaystyle\int_0^a f^3(x)  \mathrm{d}x$.
\end{enumerate}

\subsection*{三、(本题14分)}
\begin{enumerate}[start=7]
    \item 设某物体所在的空间区域为
          \[
              \Omega: x^2 + y^2 + 2z^2 \leq x + y + 2z,
          \]
          密度函数为 $x^2 + y^2 + z^2$,  求该物体的质量 $M = \displaystyle\iiint\limits_{\Omega} (x^2 + y^2 + z^2)  \mathrm{d}x  \mathrm{d}y  \mathrm{d}z$.
\end{enumerate}

\subsection*{四、(本题14分)}
\begin{enumerate}[start=8]
    \item 设函数 $f(x)$ 在闭区间 $[0, 1]$ 上具有连续导数,  $f(0) = 0,  f(1) = 1$.证明：
          \[
              \lim_{n \to \infty} n \left( \int_0^1 f(x)  \mathrm{d}x - \frac{1}{n} \sum_{k=1}^n f\left( \frac{k}{n} \right) \right) = -\frac{1}{2}.
          \]
\end{enumerate}

\subsection*{五、(本题14分)}
\begin{enumerate}[start=9]
    \item 设函数 $f(x)$ 在区间 $[0, 1]$ 上连续,  且 $I = \displaystyle\int_0^1 f(x)  \mathrm{d}x \neq 0$.证明在 $(0, 1)$ 内存在不同的两点 $x_1,  x_2$,  使得
          \[
              \frac{1}{f(x_1)} + \frac{1}{f(x_2)} = \frac{2}{I}.
          \]
\end{enumerate}

\subsection*{六、(本题14分)}
\begin{enumerate}[start=10]
    \item 设 $f(x)$ 在 $(-\infty,  +\infty)$ 上可导,  且满足
          \[
              f(x) = f(x + 2) = f(x + \sqrt{3}),
          \]
          用 Fourier 级数理论证明 $f(x)$ 为常数.
\end{enumerate}


\newpage
\section{第九届}

\subsection*{一、填空题(本题42 分,  每小题7分,  共6小题)}
\begin{enumerate}
    \item 已知可导函数 $f(x)$ 满足 $f(x) \cos x + 2 \displaystyle\int_0^x f(t) \sin t  \mathrm{d}t = x + 1$,  则 $f(x) = \underline{\quad\quad}$.
    \item 极限 $\displaystyle\lim_{n \to \infty} \sin^2 \left( \pi \sqrt{n^2 + n} \right) = \underline{\quad\quad}$.
    \item 设 $w = f(u,  v)$ 具有二阶连续偏导数,  且 $u = x - cy,  v = x + cy$,  其中 $c$ 为非零常数,  则 $w_{xx} - \dfrac{1}{c^2} w_{yy} = \underline{\quad\quad}$.
    \item 设 $f(x)$ 有二阶连续导数,  且 $f(0) = f'(0) = 0,  f''(0) = 6$,  则 $\displaystyle\lim_{x \to 0} \frac{f \left( \sin^2 x \right)}{x^4} = \underline{\quad\quad}$.
    \item 不定积分 $I = \displaystyle\int \frac{e^{-\sin x} \sin 2x}{(1 - \sin x)^2}  \mathrm{d}x = \underline{\quad\quad}$.
    \item 记由曲面 $z^2 = x^2 + y^2$ 和 $z = \sqrt{4 - x^2 - y^2}$ 围成的空间区域为 $V$,  则三重积分 $\displaystyle\iiint\limits_{V} z  \mathrm{d}x  \mathrm{d}y  \mathrm{d}z = \underline{\quad\quad}$.
\end{enumerate}

\subsection*{二、(本题14分)}
\begin{enumerate}[start=7]
    \item 设二元函数 $f(x,  y)$ 在平面上有连续的二阶偏导数.对任何角度 $\alpha$,  定义一元函数
          \[
              g_\alpha(t) = f(t \cos \alpha,  t \sin \alpha).
          \]
          若对任何 $\alpha$ 都有 $\dfrac{\mathrm{d} g_\alpha(0)}{\mathrm{d} t} = 0$ 且 $\dfrac{\mathrm{d}^2 g_\alpha(0)}{\mathrm{d} t^2} > 0$,  证明 $f(0, 0)$ 是 $f(x,  y)$ 的极小值.
\end{enumerate}

\subsection*{三、(本题14分)}
\begin{enumerate}[start=8]
    \item 设曲线 $\Gamma$ 为曲线
          \[
              \begin{cases}
                  x^2 + y^2 + z^2 = 1, & x \geqslant 0,  y \geqslant 0,  z \geqslant 0, \\
                  x + z = 1,
              \end{cases}
          \]
          上从点 $A(1, 0, 0)$ 到点 $B(0, 0, 1)$ 的一段.求曲线积分
          \[
              I = \int_{\Gamma} y  \mathrm{d}x + z  \mathrm{d}y + x  \mathrm{d}z.
          \]
\end{enumerate}

\subsection*{四、(本题15分)}
\begin{enumerate}[start=9]
    \item 设函数 $f(x) > 0$ 且在实轴上连续,  若对任意实数 $t$,  有
          $\displaystyle\int_{-\infty}^{+\infty} e^{-|t-x|} f(x)  \mathrm{d}x \leqslant 1,  $
          证明：$\forall a,  b (a < b)$,  有 $\displaystyle\int_a^b f(x)  \mathrm{d}x \leqslant \frac{b - a + 2}{2}$.
\end{enumerate}

\subsection*{五、(本题15分)}
\begin{enumerate}[start=10]
    \item 设 $\{a_n\}$ 为一个数列,  $p$ 为固定的正整数,  且 $\displaystyle\lim_{n \to \infty} (a_{n+p} - a_n) = \lambda$,  其中 $\lambda$ 为常数.证明：$\displaystyle\lim_{n \to \infty} \frac{a_n}{n} = \frac{\lambda}{p}$.
\end{enumerate}


\newpage
\section{第十届}

\subsection*{一、填空题(本题24分,  每小题6分,  共4小题)}
\begin{enumerate}
    \item 设 $\alpha \in (0, 1)$,  则 $\displaystyle\lim_{n \to \infty} \left( (n+1)^\alpha - n^\alpha \right) = \underline{\quad\quad}$.
    \item 若曲线 $y = y(x)$ 由 $
              \begin{cases}
                  x = t + \cos t, \\
                  e^y + ty + \sin t = 1
              \end{cases}
          $ 确定,  则此曲线在 $t = 0$ 对应点处的切线方程为 \underline{\quad\quad}.
    \item 积分 $\displaystyle\int \frac{\ln \left( x + \sqrt{1 + x^2} \right)}{(1 + x^2)^{3/2}}  \mathrm{d}x = \underline{\quad\quad}$.
    \item 极限 $\displaystyle\lim_{x \to 0} \frac{1 - \cos x \sqrt{\cos 2x} \sqrt[3]{\cos 3x}}{x^2} = \underline{\quad\quad}$.
\end{enumerate}

\subsection*{二、(本题8分)}
\begin{enumerate}[start=5]
    \item 设函数 $f(t)$ 在 $t \neq 0$ 时一阶连续可导,  且 $f(1) = 0$,  求函数 $f(x^2 - y^2)$,  使得曲线积分
          \[
              \int_L y \left[ 2 - f(x^2 - y^2) \right]  \mathrm{d}x + x f(x^2 - y^2)  \mathrm{d}y
          \]
          与路径无关,  其中 $L$ 为任一不与直线 $y = \pm x$ 相交的分段光滑闭曲线.
\end{enumerate}

\subsection*{三、(本题14分)}
\begin{enumerate}[start=6]
    \item 设 $f(x)$ 在区间 $[0, 1]$ 上连续,  且 $1 \leqslant f(x) \leqslant 3$.证明：
          \[
              1 \leqslant \int_0^1 f(x)  \mathrm{d}x \int_0^1 \frac{1}{f(x)}  \mathrm{d}x \leqslant \frac{4}{3}.
          \]
\end{enumerate}

\subsection*{四、(本题12分)}
\begin{enumerate}[start=7]
    \item 计算三重积分 $\displaystyle\iiint\limits_{V} (x^2 + y^2)  \mathrm{d}V$,  其中 $V$ 是由 $x^2 + y^2 + (z-2)^2 \geqslant 4,  x^2 + y^2 + (z-1)^2 \leqslant 9,  z \geqslant 0$ 所围成的空心立体.
\end{enumerate}

\subsection*{五、(本题14分)}
\begin{enumerate}[start=8]
    \item 设 $f(x,  y)$ 在区域 $D$ 内可微,  $\sqrt{\left( \dfrac{\partial f}{\partial x} \right)^2 + \left( \dfrac{\partial f}{\partial y} \right)^2} \leqslant M,  A(x_1,  y_1),  B(x_2,  y_2)$ 是 $D$ 内两点,  线段 $AB$ 包含在 $D$ 内.证明：
          \[
              |f(x_2,  y_2) - f(x_1,  y_1)| \leqslant M |AB|,
          \]
          其中 $|AB|$ 表示线段 $AB$ 的长度.
\end{enumerate}

\subsection*{六、(本题14分)}
\begin{enumerate}[start=9]
    \item 证明：对于连续函数 $f(x) > 0$,  有
          $\displaystyle \ln \int_0^1 f(x)  \mathrm{d}x \geqslant \int_0^1 \ln f(x)  \mathrm{d}x.$
\end{enumerate}

\subsection*{七、(本题14分)}
\begin{enumerate}[start=10]
    \item 已知 $\{a_k\},  \{b_k\}$ 是正项数列,  且 $b_{k+1} - b_k \geqslant \delta > 0 (k = 1,  2,  \cdots)$,  其中 $\delta$ 为一常数.证明：若级数 $\displaystyle\sum_{k=1}^{\infty} a_k$ 收敛,  则级数 $\displaystyle\sum_{k=1}^{\infty} \frac{k \sqrt[k]{(a_1 a_2 \cdots a_k)(b_1 b_2 \cdots b_k)}}{b_{k+1} b_k}$ 收敛.
\end{enumerate}


\newpage
\section{第十一届}

\subsection*{一、填空题(本题30分,  每小题6分,  共5小题)}
\begin{enumerate}
    \item (1) 极限 $\displaystyle\lim_{x \to 0} \frac{\ln(e^{\sin x} + \sqrt[3]{1 - \cos x}) - \sin x}{\arctan(4 \sqrt[3]{1 - \cos x})} = \underline{\quad\quad}$.
    \item 设隐函数 $y = y(x)$ 由方程 $y^2(x - y) = x^2$ 所确定, 则 $\displaystyle\int \frac{\mathrm{d}x}{y^2} = \underline{\quad\quad}$.
    \item 定积分 $\displaystyle\int_0^{\frac{\pi}{2}} \frac{e^x (1 + \sin x)}{1 + \cos x}  \mathrm{d}x = \underline{\quad\quad}$.
    \item 已知 $\mathrm{d}u(x,  y) = \dfrac{y  \mathrm{d}x - x  \mathrm{d}y}{3x^2 - 2xy + 3y^2}$, 则 $u(x,  y) = \underline{\quad\quad}$.
    \item 设 $a,  b,  c,  \mu > 0$, 曲面 $xyz = \mu$ 与曲面 $\dfrac{x^2}{a^2} + \dfrac{y^2}{b^2} + \dfrac{z^2}{c^2} = 1$ 相切, 则 $\mu = \underline{\quad\quad}$.
\end{enumerate}

\subsection*{二、(本题14分)}
\begin{enumerate}[start=6]
    \item 计算三重积分：$\displaystyle\iiint\limits_{\Omega} \frac{xyz}{x^2 + y^2}  \mathrm{d}x  \mathrm{d}y  \mathrm{d}z$, 其中 $\Omega$ 是由曲面 $(x^2 + y^2 + z^2)^2 = 2xy$ 围成的区域在第一卦限部分.
\end{enumerate}

\subsection*{三、(本题14分)}
\begin{enumerate}[start=7]
    \item 设 $f(x)$ 在 $[0,  +\infty)$ 上可微, $f(0) = 0$, 且存在常数 $A > 0$, 使得 $|f'(x)| \leqslant A |f(x)|$ 在 $[0,  +\infty)$ 上成立, 试证明：在 $(0,  +\infty)$ 上有 $f(x) \equiv 0$.
\end{enumerate}

\subsection*{四、(本题14分)}
\begin{enumerate}[start=8]
    \item 计算积分 $I = \displaystyle\int_0^{2\pi} \mathrm{d}\varphi \int_0^\pi e^{\sin \theta (\cos \varphi - \sin \varphi)} \sin \theta  \mathrm{d}\theta$.
\end{enumerate}

\subsection*{五、(本题14分)}
\begin{enumerate}[start=9]
    \item 设 $f(x)$ 是仅有正实根的多项式函数, 满足 $\dfrac{f'(x)}{f(x)} = - \displaystyle\sum_{n=0}^{\infty} c_n x^n$.试证：$c_n > 0  (n \geqslant 0)$, 极限 $\displaystyle\lim_{n \to \infty} \frac{1}{\sqrt[n]{c_n}}$ 存在, 且等于 $f(x)$ 的最小根.
\end{enumerate}

\subsection*{六、(本题14分)}
\begin{enumerate}[start=10]
    \item 设函数 $f(x)$ 在 $[0,  +\infty)$ 上具有连续导数, 满足
          \[
              3[3 + f^2(x)] f'(x) = 2[1 + f^2(x)]^2 e^{-x^2},
          \]
          且 $f(0) \leqslant 1$.证明：存在常数 $M > 0$, 使得当 $x \in [0,  +\infty)$ 时, 恒有 $|f(x)| \leqslant M$.
\end{enumerate}


\newpage
\section{第十二届}

\subsection*{一、填空题(本题 30 分,  每小题 6分,  共5 小题)}
\begin{enumerate}
    \item 极限 $\displaystyle\lim_{x \to 0} \frac{(x - \sin x) e^{-x^2}}{\sqrt{1 - x^3} - 1} = \underline{\quad\quad}$.
    \item 设函数 $f(x) = (x + 1)^n e^{-x^2}$, 则 $f^{(n)}(-1) = \underline{\quad\quad}$.
    \item 设 $y = f(x)$ 是由方程 $\arctan \dfrac{x}{y} = \ln \sqrt{x^2 + y^2} - \dfrac{1}{2} \ln 2 + \dfrac{\pi}{4}$ 确定的隐函数, 且满足 $f(1) = 1$, 则曲线 $y = f(x)$ 在点 $(1,  1)$ 处的切线方程为 \underline{\quad\quad}.
    \item 已知 $\displaystyle\int_0^{+\infty} \frac{\sin x}{x}  \mathrm{d}x = \frac{\pi}{2}$, 则 $\displaystyle\int_0^{+\infty} \int_0^{+\infty} \frac{\sin x \sin (x + y)}{x(x + y)}  \mathrm{d}x  \mathrm{d}y = \underline{\quad\quad}$.
    \item 设 $f(x),  g(x)$ 在 $x = 0$ 的某一邻域 $U$ 内有定义, 对任意 $x \in U,  f(x) \neq g(x)$, 且 $\displaystyle\lim_{x \to 0} f(x) = \lim_{x \to 0} g(x) = a > 0$, 则 $\displaystyle\lim_{x \to 0} \frac{[f(x)]^{g(x)} - [g(x)]^{g(x)}}{f(x) - g(x)} = \underline{\quad\quad}$.
\end{enumerate}

\subsection*{二、(本题 10 分)}
\begin{enumerate}[start=6]
    \item 设数列 $\{a_n\}$ 满足 $a_0 = 1$, 且 $a_{n+1} = \dfrac{a_n}{(n+1)(a_n+1)}$, $n \geqslant 1$.
          求极限 $\displaystyle\lim_{n \to \infty} n! a_n$.
\end{enumerate}

\subsection*{三、(本题 10 分)}
\begin{enumerate}[start=7]
    \item 设 $f(x)$ 在 $[0, 1]$ 上连续, $f(x)$ 在 $(0, 1)$ 内可导, 且 $f(0) = 0$, $f(1) = 1$.证明：
          \begin{enumerate}
              \item [(1)] 存在 $x_0 \in (0, 1)$ 使得 $f(x_0) = 2 - 3x_0$；
              \item [(2)] 存在 $\xi,  \eta \in (0, 1)$, 且 $\xi \neq \eta$, 使得 $[1 + f'(\xi)][1 + f'(\eta)] = 4$.
          \end{enumerate}
\end{enumerate}

\subsection*{四、(本题 12 分)}
\begin{enumerate}[start=8]
    \item 已知 $z = x f\left(\dfrac{y}{x}\right) + 2y\phi\left(\dfrac{x}{y}\right)$, 其中 $f,  \phi$ 均为二次可微函数.
          \begin{enumerate}
              \item [(1)] 求 $\dfrac{\partial z}{\partial x},  \dfrac{\partial^2 z}{\partial x \partial y}$；
              \item [(2)] 当 $f = \phi$, 且 $\left.\dfrac{\partial^2 z}{\partial x \partial y}\right|_{x=a} = -by^2$ 时, 求 $f(y)$.
          \end{enumerate}
\end{enumerate}

\subsection*{五、(本题 12 分)}
\begin{enumerate}[start=9]
    \item 计算 $I = \displaystyle\oint_{\Gamma} |\sqrt{3}y - x|  \mathrm{d}x - 5z  \mathrm{d}z$, 曲线 $\Gamma: \left\{\begin{array}{l}x^2 + y^2 + z^2 = 8,  \\ x^2 + y^2 = 2z\end{array}\right.$ 从 $z$ 轴正向往坐标原点看去取逆时针方向.
\end{enumerate}

\subsection*{六、(本题 12 分)}
\begin{enumerate}[start=10]
    \item 证明：$f(n) = \displaystyle\sum_{m=1}^{n} \int_{0}^{m} \cos \frac{2\pi n[x+1]}{m} \mathrm{d}x$ 等于 $n$ 的所有因子（包括 $1$ 和 $n$ 本身）之和, 其中 $[x+1]$ 表示不超过 $x+1$ 的最大整数, 并计算 $f(2021)$.
\end{enumerate}

\subsection*{七、(本题 14 分)}
\begin{enumerate}[start=11]
    \item 设 $u_n = \displaystyle\int_{0}^{1} \frac{\mathrm{d}t}{(1+t^4)^n} \ (n \geqslant 1)$.
          \begin{enumerate}
              \item[(1)] 证明数列 $\{u_n\}$ 收敛, 并求极限 $\displaystyle\lim_{n \to \infty} u_n$；
              \item[(2)] 证明级数 $\displaystyle\sum_{n=1}^{\infty} (-1)^n u_n$ 条件收敛；
              \item[(3)] 证明当 $p \geqslant 1$ 时级数$\displaystyle\sum_{n=1}^{\infty} \frac{u_n}{n^p}$ 收敛, 并求级数 $\displaystyle\sum_{n=1}^{\infty} \frac{u_n}{n}$ 的和.
          \end{enumerate}
\end{enumerate}


\newpage
\section{第十三届}

\subsection*{一、填空题(本题 30 分,  每小题 6分,  共5 小题)}
\begin{enumerate}
    \item 极限$\displaystyle\lim_{x \to +\infty} \sqrt{x^2 + x + 1} \frac{x - \ln(e^x + x)}{x} = \underline{\quad \quad} $.
    \item 设 $z = z(x,  y)$ 是由方程 $2 \sin(x + 2y - 3z) = x + 2y - 3z$ 所确定的二元隐函数, 则 $ \dfrac{\partial z}{\partial x} + \dfrac{\partial z}{\partial y} = \underline{\quad \quad} $.
    \item 设函数 $f(x)$ 连续, 且 $f(0) \neq 0$, 则 $ \displaystyle\lim_{x \to 0} \frac{2 \displaystyle\int_0^x (x - t) f(t) \mathrm{d}t}{x \displaystyle\int_0^x f(x - t) \mathrm{d}t} = \underline{\quad \quad} $.
    \item 过三条直线 $L_1 : \begin{cases} x = 0,  \\ y - z = 2,  \end{cases}$ $L_2 : \begin{cases} x = 0,  \\ x + y - z + 2 = 0 \end{cases}$ 与 $L_3 : \begin{cases} x = \sqrt{2},  \\ y - z = 0 \end{cases}$ 的圆柱面方程为 \underline{\quad \quad}.
    \item 记 $D = \{(x,  y) | x^2 + y^2 \leq \pi \}$, 则 $ \displaystyle\iint\limits_D \left( \sin x^2 \cos y^2 + x \sqrt{x^2 + y^2} \right) \mathrm{d}x\mathrm{d}y = \underline{\quad \quad} $.
\end{enumerate}

\subsection*{二、(本题 14 分)}
\begin{enumerate}[start=6]
    \item 设 $x_1 = 2021$, $x_n^2 - 2(x_n + 1)x_{n+1} + 2021 = 0 \ (n \geq 1)$, 证明数列 $\{x_n\}$ 收敛, 并求极限 $ \displaystyle\lim_{n \to \infty} x_n $.
\end{enumerate}

\subsection*{三、(本题 14 分)}
\begin{enumerate}[start=7]
    \item 设 $f(x)$ 在 $[0, +\infty)$ 上是有界连续函数, 证明：方程 $y'' + 14y' + 13y = f(x)$ 的每一个解在 $[0, +\infty)$ 上都是有界函数.
\end{enumerate}

\subsection*{四、(本题 14 分)}
\begin{enumerate}[start=8]
    \item 对于 4 次齐次函数
          \[ f(x, y, z) = a_1x^4 + a_2y^4 + a_3z^4 + 3a_4x^2y^2 + 3a_5y^2z^2 + 3a_6x^2z^2,  \]
          计算曲面积分 $ \displaystyle\varoiint_{\Sigma} f(x, y, z) \mathrm{d}S $, 其中 $\Sigma : x^2 + y^2 + z^2 = 1$.
\end{enumerate}

\subsection*{五、(本题 14 分)}
\begin{enumerate}[start=9]
    \item 设函数 $f(x)$ 在闭区间 $[a, b]$ 上有连续的二阶导数, 证明：
          \[ \lim_{n \to \infty} n^2 \left[ \int_a^b f(x) \mathrm{d}x - \frac{b-a}{n} \sum_{k=1}^n f\left( a + \frac{2k-1}{2n}(b-a) \right) \right] = \frac{(b-a)^2}{24} [f'(b) - f'(a)] .\]
\end{enumerate}

\subsection*{六、(本题 14 分)}
\begin{enumerate}[start=10]
    \item 设 $\{a_n\}$ 与 $\{b_n\}$ 均为正实数列, 满足 $a_1 = b_1 = 1$, 且 $b_n = a_n b_{n-1} - 2,  n = 2,  3,  \cdots$.又设 $\{b_n\}$ 为有界数列, 证明级数 $ \displaystyle\sum_{n=1}^{\infty} \frac{1}{a_1 a_2 \cdots a_n} $ 收敛, 并求该级数的和.
\end{enumerate}


\newpage
\section{第十三届补赛}

\subsection*{一、填空题(本题 30 分,  每小题 6分,  共5 小题)}
\begin{enumerate}
    \item 设 $x_0 = 1$, $x_n = \ln (1 + x_{n-1})(n \geq 1)$, 则 $\displaystyle\lim_{n \to \infty} nx_n = $ \underline{\quad\quad}.
    \item 积分 $I = \displaystyle\int_0^{\frac{\pi}{2}} \frac{\cos x}{1 + \tan x} \mathrm{d}x = $ \underline{\quad\quad}.
    \item 已知直线 $L : \begin{cases}
                  2x - 4y + z = 0, \\
                  3x - y - 2z = 9
              \end{cases}$ 和平面 $\Pi : 4x - y + z = 1$, 则直线 $L$ 在平面 $\Pi$ 上的投影直线方程为 \underline{\quad\quad}.
    \item $ \displaystyle\sum_{n=1}^{\infty} \arctan \frac{2}{4n^2 + 4n + 1} = $ \underline{\quad\quad}.
    \item 微分方程 $ \begin{cases}
                  (x + 1) \dfrac{\mathrm{d}y}{\mathrm{d}x} + 1 = 2e^{-y}, \\
                  y(0) = 0
              \end{cases} $ 的解为 \underline{\quad\quad}.
\end{enumerate}

\subsection*{二、(本题 14 分)}
\begin{enumerate}[start=6]
    \item 设 $f(x) = -\dfrac{1}{2} \left( 1 + \dfrac{1}{e} \right) + \displaystyle\int_{-1}^{1} |x - t| e^{-t^2} \mathrm{d}t$, 证明：$f(x)$ 在区间 $(-1,  1)$ 内有且仅有两个实根.
\end{enumerate}

\subsection*{三、(本题 14 分)}
\begin{enumerate}[start=7]
    \item 设函数 $f(x,  y)$ 在闭区域 $D = \{(x,  y)| x^2 + y^2 \leq 1\}$ 上具有二阶连续偏导数, 且满足 $ \dfrac{\partial^2 f}{\partial x^2} + \dfrac{\partial^2 f}{\partial y^2} = x^2 + y^2 $, 求极限 $ \displaystyle\lim_{r \to 0+} \frac{\displaystyle\iint\limits_{x^2 + y^2 \leq r^2} \left( x \frac{\partial f}{\partial x} + y \frac{\partial f}{\partial y} \right) \mathrm{d}x\mathrm{d}y}{(\tan r - \sin r)^2} $.
\end{enumerate}

\subsection*{四、(本题 14 分)}
\begin{enumerate}[start=8]
    \item 设对于 $\mathbb{R}^3$ 中半空间 $\{(x,  y,  z) \in \mathbb{R}^3 | x > 0\}$ 内任意有向光滑封闭曲面 $S$, 都有 \[ \varoiint_{S} xf'(x)\mathrm{d}y\mathrm{d}z + y[xf(x) - f'(x)]\mathrm{d}z\mathrm{d}x - xz[\sin x + f'(x)]\mathrm{d}x\mathrm{d}y = 0 , \]其中 $f(x)$ 在 $(0,  +\infty)$ 上二阶导数连续, 且 $ \displaystyle\lim_{x \to 0+} f(x) = \displaystyle\lim_{x \to 0+} f'(x) = 0 $, 求 $f(x)$.
\end{enumerate}

\subsection*{五、(本题 14 分)}
\begin{enumerate}[start=9]
    \item 设函数 $f(x) = \displaystyle\int_0^x \left(1 - \frac{[u]}{u}\right) \mathrm{d}u$, 其中 $[u]$ 表示不超过 $u$ 的最大整数, 试讨论 $\displaystyle \int_1^{+\infty} \frac{e^{f(x)}}{x^p} \cos \left(x^2 - \frac{1}{x^2}\right) \mathrm{d}x $ 的敛散性, 其中 $p > 0$.
\end{enumerate}

\subsection*{六、(本题 14 分)}
\begin{enumerate}[start=10]
    \item 设正数列 $\{a_n\}$ 单调减少趋于零, $f(x) = \displaystyle\sum_{n=1}^\infty a_n^n x^n$, 证明：若级数 $\sum_{n=1}^\infty a_n$ 发散, 则积分 $\displaystyle\int_1^{+\infty} \frac{\ln f(x)}{x^2} \mathrm{d}x $ 也发散.
\end{enumerate}


\newpage
\section{第十四届}

\subsection*{一、填空题(本题 30 分,  每小题 6分,  共5 小题)}
\begin{enumerate}
    \item 极限 $\displaystyle\lim_{x \to 0} \frac{1 - \sqrt{1 - x^2} \cos x}{1 + x^2 - \cos^2 x} = \underline{\quad\quad}. $
    \item 设 $f(x) = \begin{cases}
                  1, & x > 0,    \\
                  0, & x \leq 0,
              \end{cases}$ $g(x) = \begin{cases}
                  x - 1, & x \geq 1, \\
                  1 - x, & x < 1,
              \end{cases}$ 则复合函数 $f[g(x)]$ 的间断点为 $x = \underline{\quad\quad}.$
    \item 极限 $\displaystyle \lim_{x \to 1^-} (1 - x)^3 \sum_{n=1}^\infty n^2 x^n = \underline{\quad\quad}. $
    \item 微分方程 $ \dfrac{\mathrm{d}y}{\mathrm{d}x} x \ln x \cdot \sin y + \cos y (1 - x \cos y) = 0 $ 的通解为 \underline{\quad\quad}.
    \item 记 $D = \{(x,  y) | 0 \leq x + y \leq \dfrac{\pi}{2},  0 \leq x - y \leq \dfrac{\pi}{2}\}$,  则 \[ \iint\limits_D y \sin(x + y) \mathrm{d}x \mathrm{d}y = \underline{\quad\quad}. \]
\end{enumerate}

\subsection*{二、(本题 14 分)}
\begin{enumerate}[start=6]
    \item 记向量 $\overrightarrow{OA}$ 与 $\overrightarrow{OB}$ 的夹角为 $\alpha$, $|\overrightarrow{OA}| = 1$, $|\overrightarrow{OB}| = 2$, $\overrightarrow{OP} = (1 - \lambda) \overrightarrow{OA}$, $\overrightarrow{OQ} = \lambda \overrightarrow{OB}$, 其中 $0 \leq \lambda \leq 1$.
          \begin{enumerate}
              \item [(1)] 问当 $\lambda$ 为何值时, $|\overrightarrow{PQ}|$ 取得最小值；
              \item [(2)] 设 (1) 中的 $\lambda$ 满足 $0 < \lambda < \dfrac{1}{5}$, 求夹角 $\alpha$ 的取值范围.
          \end{enumerate}
\end{enumerate}

\subsection*{三、(本题 14 分)}
\begin{enumerate}[start=7]
    \item 设函数 $f(x)$ 在 $(-1,  1)$ 上二阶可导, $f(0) = 1$, 且当 $x \geq 0$ 时, $f(x) \geq 0$, $f'(x) \leq 0$, $f''(x) \leq f(x)$, 证明：$f'(0) \geq -\sqrt{2}$.
\end{enumerate}

\subsection*{四、(本题 14 分)}
\begin{enumerate}[start=8]
    \item 证明：对任意正整数 $n$, 恒有
          \[ \int_{0}^{\frac{\pi}{2}} x \left( \frac{\sin nx}{\sin x} \right)^{4} \mathrm{d}x \leq \left( \frac{n^2}{4} - \frac{1}{8} \right) \pi^{2} .\]
\end{enumerate}

\subsection*{五、(本题 14 分)}
\begin{enumerate}[start=9]
    \item 设 $z = f(x,  y)$ 是区域 $D = \{(x,  y) | 0 \leq x \leq 1,  0 \leq y \leq 1\}$ 上的可微函数, $f(0,  0) = 0$, 且 $\mathrm{d}z|_{(0,  0)} = 3\mathrm{d}x + 2\mathrm{d}y$, 求极限
          $\displaystyle \lim_{x \to 0^{+}} \frac{\displaystyle\int_{0}^{x^{2}} \mathrm{d}t \int_{x}^{\sqrt{t}} f(t,  u) \mathrm{d}u}{1 - \sqrt[4]{1 - x^{4}}} $.
\end{enumerate}

\subsection*{六、(本题 14 分)}
\begin{enumerate}[start=10]
    \item 设正项级数 $\displaystyle\sum_{n=1}^{\infty} a_n$ 收敛, 证明：存在收敛的正项级数 $\displaystyle\sum_{n=1}^{\infty} b_n$, 使得
          $ \displaystyle\lim_{n \to \infty} \frac{a_n}{b_n} = 0 .$
\end{enumerate}


\newpage
\section{第十四届补赛一}

\subsection*{一、填空题(本题 30 分,  每小题 6分,  共5 小题)}
\begin{enumerate}
    \item 函数 $f(x) = 2x^3 - 6x + 1$ 在区间 $\left[ \dfrac{1}{2},  2 \right]$ 上的值域为\underline{\quad\quad}.
    \item 设 $x \in (-\infty,  +\infty)$, 则极限 $\displaystyle \lim_{n \to \infty} \frac{1}{n^2} \sum_{i=1}^n \sqrt{(ne^x + i)(ne^x + i + 1)} = \underline{\quad\quad}. $
    \item 设函数 $y = y(x)$ 由方程 $x^2 + 4xy + e^y = 1$ 确定, 则 $ \left. \dfrac{\mathrm{d}^2 y}{\mathrm{d}x^2} \right|_{x=0} = \underline{\quad\quad}. $
    \item 设 $D = \{(x,  y)| x^2 + y^2 \leq 1 \}$, 则 $ \displaystyle\iint\limits_D (x + y - x^2) \mathrm{d}x\mathrm{d}y = \underline{\quad\quad}. $
    \item 设可微函数 $f(x,  y)$ 对任意 $u,  v,  t$ 满足 $f(tu,  tv) = t^2 f(u,  v)$, 点 $P(1,  -1,  2)$ 位于曲面 $z = f(x,  y)$ 上, 又设 $f'_x(1,  -1) = 3$, 则该曲面在点 $P$ 处的切平面方程为\underline{\quad\quad}..
\end{enumerate}

\subsection*{二、(本题 14 分)}
\begin{enumerate}[start=6]
    \item 设函数 $z = f(u)$ 在区间 $(0,  +\infty)$ 上具有二阶连续导数, $u = \sqrt{x^2 + y^2}$, 且满足 $\dfrac{\partial^2 z}{\partial x^2} + \dfrac{\partial^2 z}{\partial y^2} = x^2 + y^2$, 求函数 $f(u)$ 的表达式.
\end{enumerate}

\subsection*{三、(本题 14 分)}
\begin{enumerate}[start=7]
    \item 设曲线 $C: x^3 + y^3 - \dfrac{3}{2}xy = 0$.
          \begin{enumerate}
              \item [(1)] 已知曲线 $C$ 存在斜渐近线, 求其斜渐近线的方程；
              \item [(2)] 求由曲线 $C$ 所围成的平面图形的面积.
          \end{enumerate}
\end{enumerate}

\subsection*{四、(本题 14 分)}
\begin{enumerate}[start=8]
    \item 证明：当 $\alpha > 0$ 时, $\left(\dfrac{2\alpha + 2}{2\alpha + 1}\right)^{\alpha+1} > \left(\dfrac{2\alpha + 1}{2\alpha}\right)^{\alpha}$.
\end{enumerate}

\subsection*{五、(本题 14 分)}
\begin{enumerate}[start=9]
    \item 计算曲线积分 $I = \displaystyle\oint\limits_{\Gamma}(y^2 + z^2)\mathrm{d}x + (z^2 + x^2)\mathrm{d}y + (x^2 + y^2)\mathrm{d}z$, 其中 $\Gamma : \begin{cases}
                  x^2 + y^2 + z^2 = 2Rx, \\
                  x^2 + y^2 = 2rx
              \end{cases} (0 < r < R,  z \geq 0)$, 方向与 $z$ 轴正向符合右手螺旋法则.
\end{enumerate}

\subsection*{六、(本题 14 分)}
\begin{enumerate}[start=10]
    \item 证明方程 $x = \tan\sqrt{x}$ 有无穷多个正根, 且所有正根 $\{r_n\}$ 可以按递增顺序排列为 $0 < r_1 < r_2 < \cdots < r_n < \cdots$, 并讨论级数 $\displaystyle\sum_{n=1}^{\infty}(\cot\sqrt{r_n})^\lambda$ 的收敛性, 其中 $\lambda$ 是正常数.
\end{enumerate}


\newpage
\section{第十四届补赛二}

\subsection*{一、填空题(本题 30 分,  每小题 6分,  共5 小题)}
\begin{enumerate}
    \item 极限 $ \displaystyle\lim_{n \to \infty} \frac{1}{n^3} \left[ 1^2 + 3^2 + \cdots + (2n - 1)^2 \right] = \underline{\quad\quad}. $
    \item 设函数 $f(x)$ 在 $x = 1$ 的某一邻域内可微, 且满足
          \[ f(1 + x) - 3f(1 - x) = 4 + 2x + o(x),  \]
          其中 $o(x)$ 是当 $x \to 0$ 时 $x$ 的高阶无穷小, 则曲线 $y = f(x)$ 在点 $(1,  f(1))$ 处的切线方程为 \underline{\quad\quad}.
    \item 设 $y = y(x)$ 是初值问题
          $ \begin{cases}
                  y'' - 2y' - 3y = 1, \\
                  y(0) = 0,  y'(0) = 1
              \end{cases} $
          的解, 则 $y(x) = \underline{\quad\quad}.$
    \item 设可微函数 $z = z(x,  y)$ 满足
          $ x^2 \dfrac{\partial z}{\partial x} + y^2 \dfrac{\partial z}{\partial y} = 2z^2,  $
          又设 $u = x$,  $v = \dfrac{1}{y} - \dfrac{1}{x}$,  $w = \dfrac{1}{z} - \dfrac{1}{x}$, 则对函数 $w = w(u,  v)$, 偏导数
          $ \dfrac{\partial w}{\partial u} \bigg|_{\substack{u=2 \\ v=1}} = \underline{\quad\quad}. $
    \item 设 $a > 0$, 则均匀曲面 $x^2 + y^2 + z^2 = a^2$ ($x \geq 0$,  $y \geq 0$,  $z \geq 0$) 的重心坐标为 \underline{\quad\quad}.
\end{enumerate}

\subsection*{二、(本题 14 分)}
\begin{enumerate}[start=6]
    \item 设函数 $f(x) = e^{-x} \displaystyle\int_{0}^{x} \frac{t^{2023}}{1+t^2} \mathrm{d}t$, 正整数 $n \leq 2023$, 求导数 $f^{(n)}(0)$.
\end{enumerate}

\subsection*{三、(本题 14 分)}
\begin{enumerate}[start=7]
    \item 设函数 $f(x)$ 在区间 $(0, 1)$ 内有定义,
          $ \displaystyle\lim_{x \to 0^+} f(x) = 0,  $
          且
          $ \displaystyle\lim_{x \to 0^+} \frac{f(x) - f\left(\dfrac{x}{3}\right)}{x} = 0. $
          证明：
          $ \displaystyle\lim_{x \to 0^+} \frac{f(x)}{x} = 0. $
\end{enumerate}

\subsection*{四、(本题 14 分)}
\begin{enumerate}[start=8]
    \item 设函数 $f(x)$ 在区间 $[0, 1]$ 上连续, 在 $(0, 1)$ 内可导, 且 $f(0) = 0$, $f(1) = 2$.证明：存在两两互异的点 $\xi_1,  \xi_2,  \xi_3 \in (0, 1)$, 使得 $f'(\xi_1)f'(\xi_2)\sqrt{1-\xi_3} \geq 2$.
\end{enumerate}

\subsection*{五、(本题 14 分)}
\begin{enumerate}[start=9]
    \item 设 $f(x)$ 是 $[-1, 1]$ 上的连续的偶函数, 计算曲线积分：
          $ I = \displaystyle\oint\limits_{L} \frac{x^2 + y^2}{2\sqrt{1-x^2}} \mathrm{d}x + f(x) \mathrm{d}y,  $
          其中曲线 $L$ 为正向圆周 $x^2 + y^2 = -2y$.
\end{enumerate}

\subsection*{六、(本题 14 分)}
\begin{enumerate}[start=10]
    \item 设函数 $f(x) = \displaystyle\int_{0}^{x} \frac{\ln(1+t)}{1+e^{-t} \sin^3 t} \mathrm{d}t$ ($x > 0$), 证明级数
          $ \displaystyle\sum_{n=1}^{\infty} f\left(\frac{1}{n}\right) $
          收敛, 且
          $ \dfrac{1}{3} < \displaystyle\sum_{n=1}^{\infty} f\left(\frac{1}{n}\right) < \dfrac{5}{6}. $
\end{enumerate}


\newpage
\section{第十五届非数A}

\subsection*{一、填空题(本题 30 分,  每小题 6分,  共5 小题)}
\begin{enumerate}
    \item 极限 $\displaystyle \lim_{x \to 3} \frac{\sqrt{x^3 + 9 - 6}}{2 - \sqrt{x^3 - 23}} = \underline{\quad\quad} .$
    \item 设函数 $z = f(x^2 - y^2, xy)$, 且 $f(u, v)$ 具有连续的二阶偏导数, 则 $ \dfrac{\partial^2 z}{\partial x \partial y} = \underline{\quad\quad} .$
    \item 设 $n$ 为正整数, $f(x) = \dfrac{1}{x^2 - 3x + 2}$, 则 $f^{(n)}(0) = \underline{\quad\quad}. $
    \item 级数 $\displaystyle \sum_{n=1}^{\infty} \frac{(-1)^{n-1}}{n(2n-1)} x^{2n} $ 的收敛域为 $\underline{\quad\quad}.$
    \item 设曲面 $\Sigma$ 是平面 $y + z = 5$ 被柱面 $x^2 + y^2 = 25$ 所截得的部分, 则曲面积分 $ \displaystyle\iint\limits_{\Sigma} (x + y + z) \mathrm{d}S = \underline{\quad\quad}. $
\end{enumerate}

\subsection*{二、(本题 14 分)}
\begin{enumerate}[start=6]
    \item 解微分方程：$(x^2 + y^2 + 3) \dfrac{\mathrm{d}y}{\mathrm{d}x} = 2x \left( 2y - \dfrac{x^2}{y} \right)$.
\end{enumerate}

\subsection*{三、(本题 14 分)}
\begin{enumerate}[start=7]
    \item 设 $\Sigma_1$ 是以 $(0,4,0)$ 为顶点且与曲面 $\Sigma_2 : \dfrac{x^2}{3} + \dfrac{y^2}{4} + \dfrac{z^2}{3} = 1 (y > 0)$ 相切的圆锥面, 求曲面 $\Sigma_1$ 与 $\Sigma_2$ 所围成的空间区域的体积.
\end{enumerate}

\subsection*{四、(本题 14 分)}
\begin{enumerate}[start=8]
    \item 设 $I_n = n \displaystyle\int_{1}^{a} \frac{\mathrm{d}x}{1 + x^n}$, 其中 $a > 1$, 求极限 $\displaystyle\lim_{n \to \infty} I_n$.
\end{enumerate}

\subsection*{五、(本题 14 分)}
\begin{enumerate}[start=9]
    \item 设 $f(x)$ 在 $[0,1]$ 上有连续的导数, 且 $f(0) = 0$.求证：
          \[ \int_{0}^{1} f^2(x) \mathrm{d}x \leq 4 \int_{0}^{1} (1 - x)^2 |f'(x)|^2 \mathrm{d}x, \]
          并求使得上式成为等式的 $f(x)$.
\end{enumerate}

\subsection*{六、(本题 14 分)}
\begin{enumerate}[start=10]
    \item 设数列 $\{x_n\}$ 满足 $x_0 = \dfrac{1}{3}, x_{n+1} = \dfrac{x_n^2}{1 - x_n + x_n^2}, n \geq 0$.证明无穷级数 $\displaystyle\sum_{n=0}^{\infty} x_n$ 收敛, 并求其和.
\end{enumerate}


\newpage
\section{第十五届非数B}

\subsection*{一、填空题(本题 30 分,  每小题 6分,  共5 小题)}
\begin{enumerate}
    \item 极限 $\displaystyle \lim_{x \to \infty} \left( \frac{x + 3}{x + 2} \right)^{2x - 1} = \underline{\quad\quad}. $
    \item 设函数 $z = f(x^2 - y^2, xy)$, 且 $f(u, v)$ 具有连续的二阶偏导数, 则 $ \dfrac{\partial^2 z}{\partial x \partial y} = \underline{\quad\quad} $.
    \item 设曲线 $y = \ln(1 + ax) + 1$ 与曲线 $y = 2xy^3 + b$ 在点 $(0, 1)$ 处相切, 则 $a + b = \underline{\quad\quad} $.
    \item 设函数 $y = y(x)$ 由方程 $y = 1 + \arctan(xy)$ 确定, 则 $y'(0) = \underline{\quad\quad} $.
    \item 计算积分：$\displaystyle \int_0^1 \mathrm{d}x \int_x^{\sqrt{x}} \frac{\cos y}{y} \mathrm{d}y = \underline{\quad\quad} $.
\end{enumerate}

\subsection*{二、(本题 14 分)}
\begin{enumerate}[start=6]
    \item 设曲线 $y = 3ax^2 + 2bx + \ln c$ 经过 $(0,0)$ 点, 且当 $0 \leq x \leq 1$ 时, $y \geq 0$.又该曲线与直线 $x = 1$, $x$ 轴所围成的平面图形 $D$ 的面积为 1.试求常数 $a, b, c$ 的值, 使得 $D$ 绕 $x$ 轴旋转一周所得旋转体的体积最小.
\end{enumerate}

\subsection*{三、(本题 14 分)}
\begin{enumerate}[start=7]
    \item 解微分方程：$(x^2 + y^2 + 3)\dfrac{\mathrm{d}y}{\mathrm{d}x} = 2x\left(2y - \dfrac{x^2}{y}\right)$.
\end{enumerate}

\subsection*{四、(本题 14 分)}
\begin{enumerate}[start=8]
    \item 求幂级数 $\displaystyle\sum_{n=1}^{\infty}\frac{(-1)^{n-1}}{n(2n-1)}x^{2n}$ 的收敛域与和函数.
\end{enumerate}

\subsection*{五、(本题 14 分)}
\begin{enumerate}[start=9]
    \item 设 $f(x)$ 在 $[0,1]$ 上可导, 满足 $f(0)>0$, $f(1)>0$, 且 $\displaystyle\int_{0}^{1}f(x)\mathrm{d}x=0$.
          证明：
          \begin{enumerate}
              \item [(1)] $f(x)$ 在 $(0,1)$ 内至少有两个零点；
              \item [(2)] 在 $(0,1)$ 内至少存在一点 $\xi$, 使得 $f'(\xi) + 3f^3(\xi) = 0$.
          \end{enumerate}
\end{enumerate}

\subsection*{六、(本题 14 分)}
\begin{enumerate}[start=10]
    \item 设 $f(x)$ 在 $[0,1]$ 上有连续的导数, 且 $f(0)=0$.求证：
          \[ \int_{0}^{1}f^2(x)\mathrm{d}x \leq 4\int_{0}^{1}(1-x)^2|f'(x)|^2\mathrm{d}x, \]
          并求使得上述式成为等式的 $f(x)$.
\end{enumerate}


\newpage
\section{第十六届非数A}

\subsection*{一、填空题(本题 30 分,  每小题 6分,  共5 小题)}
\begin{enumerate}
    \item $\displaystyle\int_{0}^{1} \ln(1 + x^2) \mathrm{d}x =$\underline{\quad\quad} .
    \item 设 $D: x^2 + y^2 \leqslant r^2 (r > 0)$, 则 $\displaystyle\lim_{r \to 0^+} \frac{\displaystyle\iint\limits_{D} (e^{x^2 + y^2} - 1) \mathrm{d}x \mathrm{d}y}{r^4} =$\underline{\quad\quad} .
    \item 已知 $z = f(xy, e^{x+y})$, 且 $f(x, y)$ 具有二阶连续偏导数, 则 $\dfrac{\partial^2 z}{\partial x \partial y} = $\underline{\quad\quad} .
    \item 直线 $L: \frac{x-1}{1} = \frac{y}{1} = \frac{z-1}{-1}$ 在平面 $\pi: x - y + 2z - 1 = 0$ 上的投影直线 $L_0$ 的单位方向向量为\underline{\quad\quad} .
    \item 设 $L$ 为圆周 $x^2 + y^2 = 9$, 取逆时针方向, 则第二型曲线积分
          \[
              \int_{L} \frac{-y}{4x^2 + y^2} \mathrm{d}x + \frac{x}{4x^2 + y^2} \mathrm{d}y =\underline{\quad\quad} .
          \]
\end{enumerate}

\subsection*{二、(本题 14 分)}
\begin{enumerate}[start=6]
    \item 求微分方程 $(x^3 - y^2) \mathrm{d}x + (x^2 y + xy) \mathrm{d}y = 0$ 的通解.
\end{enumerate}

\subsection*{三、(本题 14 分)}
\begin{enumerate}[start=7]
    \item 设函数
          $f(x) = \begin{cases}
                  \dfrac{1}{\ln(1 + x)} - \dfrac{1}{x}, & 0 < x \leq 1, \\[8pt]
                  k,                                    & x = 0
              \end{cases}$
          \begin{enumerate}
              \item[(1)] 求常数 $k$ 的值, 使得 $f(x)$ 在区间 $[0, 1]$ 上连续；
              \item[(2)]  对 (1) 中 $k$ 的值, 求函数 $f(x)$ 的最小值 $\lambda$ 与最大值 $\mu$.
          \end{enumerate}
\end{enumerate}

\subsection*{四、(本题 14 分)}
\begin{enumerate}[start=8]
    \item 设 $S$ 是上半球面 $x^2 + y^2 + z^2 = R^2 (z \geqslant 0)$, 方向取上侧, 计算
          \[
              I = \iiint\limits_{S} (x^2 - x) \mathrm{d}y \mathrm{d}z + (y^2 - y) \mathrm{d}z \mathrm{d}x + (z^2 - z) \mathrm{d}x \mathrm{d}y.
          \]
\end{enumerate}

\subsection*{五、(本题 14 分)}
\begin{enumerate}[start=9]
    \item 设 $f(x)$ 是定义在 $(-\infty, +\infty)$ 上具有连续导数的非负函数, 且存在 $M > 0$, 使得对任意的 $x, y \in (-\infty, +\infty)$, 有 $|f'(x) - f'(y)| \leqslant M |x - y|$.证明：对于任意实数 $x$, 恒有 $[f'(x)]^2 \leqslant 2Mf(x)$.
\end{enumerate}

\subsection*{六、(本题 14 分)}
\begin{enumerate}[start=10]
    \item 证明：级数 $\displaystyle\sum_{n=1}^{\infty} \sum_{k=1}^{\infty} \frac{(-1)^{[\sqrt{n}]}}{n^2 + k^2}$ 收敛, 其中 $[x]$ 表示不超过 $x$ 的最大整数.
\end{enumerate}


\newpage
\section{第十六届非数B}

\subsection*{一、填空题(本题 30 分,  每小题 6分,  共5 小题)}
\begin{enumerate}
    \item 极限 $\displaystyle\lim_{x \to 0} \frac{\ln(1 + x^2) - x^2}{\sin^4 x} =$\underline{\quad\quad} .
    \item $\displaystyle\int_{-2}^{2} (x^3 \cos^5 x + \sqrt{4 - x^2}) \mathrm{d}x =$\underline{\quad\quad} .
    \item 设 $y = y(x)$ 由方程 $e^{3y} + \\displaystyleint_{0}^{x+y} \cos t^2 \mathrm{d}t = 1$ 确定, 则 $\left. \dfrac{\mathrm{d}^2 y}{\mathrm{d}x^2} \right|_{(0,0)} =$\underline{\quad\quad} .
    \item 已知 $z = f(xy, e^{x+y})$, 且 $f(x, y)$ 具有二阶连续偏导数, 则 $\dfrac{\partial^2 z}{\partial x \partial y} =$ \underline{\quad\quad} .
    \item 设 $D: x^2 + y^2 \leqslant r^2 (r > 0)$, 则 $\displaystyle\lim_{r \to 0^+} \frac{\displaystyle\iint\limits_{D} (e^{x^2 + y^2} - 1) \mathrm{d}x \mathrm{d}y}{r^4} =$\underline{\quad\quad} .
\end{enumerate}

\subsection*{二、(本题 14 分)}
\begin{enumerate}[start=6]
    \item 用 Lagrange 乘子法求函数 $f(x, y, z) = x + 2y + 3z$ 在平面 $x - y + z = 1$ 和圆柱面 $x^2 + y^2 = 1$ 的交线上的最大值.
\end{enumerate}

\subsection*{三、(本题 14 分)}
\begin{enumerate}[start=7]
    \item 求函数 $f(x) = \dfrac{1}{(x - 1)(x + 3)}$ 在 $x_0 = 2$ 处的 Taylor 级数, 并确定它的收敛域.
\end{enumerate}

\subsection*{四、(本题 14 分)}
\begin{enumerate}[start=8]
    \item 求微分方程 $(x^3 - y^2) \mathrm{d}x + (x^2 y + xy) \mathrm{d}y = 0$ 的通解.
\end{enumerate}

\subsection*{五、(本题 14 分)}
\begin{enumerate}[start=9]
    \item 已知 $\displaystyle\int_{0}^{+\infty} \frac{\sin x}{x} \mathrm{d}x = \frac{\pi}{2}$.
          \begin{enumerate}
              \item[(1)] 计算 $I_1 = \displaystyle\int_{0}^{+\infty} \left(\frac{\sin x}{x}\right)^2 \mathrm{d}x$.
              \item[(2)] 计算 $I_2 = \displaystyle\int_{0}^{+\infty} \left(\frac{\sin x}{x}\right)^3 \mathrm{d}x$.
          \end{enumerate}
\end{enumerate}

\subsection*{六、(本题 14 分)}
\begin{enumerate}[start=10]
    \item 设 $f(x)$ 是定义在 $(-\infty, +\infty)$ 上具有连续导数的非负函数, 且存在 $M > 0$, 使得对任意的 $x, y \in (-\infty, +\infty)$, 有 $|f'(x) - f'(y)| \leqslant M |x - y|$.证明：对于任意实数 $x$, 恒有 $[f'(x)]^2 \leqslant 2Mf(x)$.
\end{enumerate}

